\textbf{Phần 1. Trắc nghiệm lựa chọn}
\setcounterpageref{ex}{0}

\begin{ex} %[2D10N54]
Số gồm $7$ trăm, $2$ chục và $0$ đơn vị là
\choice
{\True $720$}
{$702$}
{$270$}
{$72$}
\loigiai{
Số đó là $720$.\\
\textbf{Đáp án của câu hỏi là A}
}
\end{ex}

\begin{ex} %[2D10N54]
Số $805$ đọc là
\choice
{Tám trăm năm mươi}
{Tám không năm}
{\True Tám trăm linh năm}
{Tám trăm năm}
\loigiai{
Đọc là Tám trăm linh năm.\\
\textbf{Đáp án của câu hỏi là C}
}
\end{ex}

\begin{ex} %[2D10H54]
Số liền sau của số $899$ là
\choice
{\True $900$}
{$898$}
{$901$}
{$890$}
\loigiai{
$899 + 1 = 900$.\\
\textbf{Đáp án của câu hỏi là A}
}
\end{ex}

\begin{ex} %[2D10H54]
Số liền trước của số $1\,000$ là
\choice
{$1001$}
{$990$}
{$900$}
{\True $999$}
\loigiai{
$1000 - 1 = 999$.\\
\textbf{Đáp án của câu hỏi là D}
}
\end{ex}

\begin{ex} %[2D10H54]
Tính: $500 + 300 = \dots$
\choice
{$700$}
{\True $800$}
{$900$}
{$600$}
\loigiai{
$500 + 300 = 800$.\\
\textbf{Đáp án của câu hỏi là B}
}
\end{ex}

\begin{ex} %[2D10H54]
Số $642$ được viết thành tổng là
\choice
{\True $600 + 40 + 2$}
{$600 + 4 + 2$}
{$60 + 40 + 2$}
{$640 + 20$}
\loigiai{
$642 = 600 + 40 + 2$.\\
\textbf{Đáp án của câu hỏi là A}
}
\end{ex}

\begin{ex} %[2D10V54]
Điền dấu thích hợp: $789 \dots 798$.
\choice
{$>$}
{$=$}
{\True $<$}
{Không so sánh được}
\loigiai{
$789 < 798$.\\
\textbf{Đáp án của câu hỏi là C}
}
\end{ex}

\begin{ex} %[2D10V54]
Tìm số tròn chục lớn nhất có 3 chữ số:
\choice
{$900$}
{$999$}
{$1000$}
{\True $990$}
\loigiai{
Số đó là 990.\\
\textbf{Đáp án của câu hỏi là D}
}
\end{ex}

\begin{ex} %[2D10V54]
Sắp xếp các số $456, 654, 546, 465$ theo thứ tự từ bé đến lớn:
\choice
{$654, 546, 465, 456$}
{\True $456, 465, 546, 654$}
{$456, 546, 465, 654$}
{$465, 456, 546, 654$}
\loigiai{
$456 < 465 < 546 < 654$.\\
\textbf{Đáp án của câu hỏi là B}
}
\end{ex}

\begin{ex} %[2D10V54]
Tính giá trị biểu thức: $900 - 400 + 200$.
\choice
{$600$}
{$800$}
{$500$}
{\True $700$}
\loigiai{
$500 + 200 = 700$.\\
\textbf{Đáp án của câu hỏi là D}
}
\end{ex}

\begin{ex} %[2D10V54]
Số bé nhất có ba chữ số khác nhau mà chữ số hàng trăm là 5 là
\choice
{$500$}
{\True $501$}
{$510$}
{$512$}
\loigiai{
Số đó là 501.\\
\textbf{Đáp án của câu hỏi là B}
}
\end{ex}

\begin{ex} %[2D10C54]
Tìm số có 3 chữ số biết chữ số hàng trăm là 8, chữ số hàng đơn vị là 5 và chữ số hàng chục bằng hiệu của chữ số hàng trăm và hàng đơn vị.
\choice
{$853$}
{$825$}
{\True $835$}
{$845$}
\loigiai{
Hàng chục: $8 - 5 = 3$. Số đó là $835$.\\
\textbf{Đáp án của câu hỏi là C}
}
\end{ex}

\textbf{Phần 2. Trắc nghiệm đúng sai}
\setcounterpageref{ex}{0}

\begin{ex} %[2D10H54]
Xét tính đúng sai về số trong phạm vi 1 000:
\choiceTF
{\True Số 1 000 gồm 10 trăm.}
{\True $999$ là số lớn nhất có ba chữ số.}
{Số 605 đọc là Sáu trăm năm mươi.}
{\True Số 605 đọc là Sáu trăm linh năm.}
\loigiai{
\begin{itemize}
    \item 1000 = 10 trăm: Đúng.
    \item 999 là số lớn nhất có 3 chữ số: Đúng.
    \item 605 đọc là Sáu trăm linh năm. Ý 3 Sai, Ý 4 Đúng.
\end{itemize}
\begin{itemize}
    \item a) \textbf{Đúng}.
    \item b) \textbf{Đúng}.
    \item c) \textbf{Sai}.
    \item d) \textbf{Đúng}.
\end{itemize}
}
\end{ex}

\begin{ex} %[2D10V54]
Cho 4 số: $512, 521, 215, 251$. Xét tính đúng sai:
\choiceTF
{\True Số lớn nhất là 521.}
{\True Số bé nhất là 215.}
{Xếp tăng dần: $251, 215, 512, 521$.}
{\True Xếp tăng dần đúng là $215, 251, 512, 521$.}
\loigiai{
\begin{itemize}
    \item Lớn nhất 521: Đúng.
    \item Bé nhất 215: Đúng.
    \item Thứ tự đúng là $215 < 251$. Ý 3 Sai, Ý 4 Đúng.
\end{itemize}
\begin{itemize}
    \item a) \textbf{Đúng}.
    \item b) \textbf{Đúng}.
    \item c) \textbf{Sai}.
    \item d) \textbf{Đúng}.
\end{itemize}
}
\end{ex}

\textbf{Phần 3. Trả lời ngắn}
\setcounterpageref{ex}{0}

\begin{ex} %[2D10H54]
Viết số gồm 6 trăm, 3 chục và 8 đơn vị.
\shortans{638}
\loigiai{
\textbf{Đáp số:} $638$\\
638.
}
\end{ex}

\begin{ex} %[2D10H54]
Tìm số liền trước của số $500$.
\shortans{499}
\loigiai{
\textbf{Đáp số:} $499$\\
$500 - 1 = 499$.
}
\end{ex}

\begin{ex} %[2D10V54]
Tìm số liền sau của số $699$.
\shortans{700}
\loigiai{
\textbf{Đáp số:} $700$\\
$699 + 1 = 700$.
}
\end{ex}

\begin{ex} %[2D10V54]
Tính giá trị của tổng: $700 + 40 + 9$.
\shortans{749}
\loigiai{
\textbf{Đáp số:} $749$\\
749.
}
\end{ex}

\begin{ex} %[2D10V54]
Tính: $800 - 500 + 300$.
\shortans{600}
\loigiai{
\textbf{Đáp số:} $600$\\
$300 + 300 = 600$.
}
\end{ex}

\begin{ex} %[2D10C54]
Tìm số tròn chục bé nhất có ba chữ số khác nhau ($120$).
\shortans{120}
\loigiai{
\textbf{Đáp số:} $120$\\
120.
}
\end{ex}
